\chapter{Implementation and Evaluation}
\label{ch:impl}

\section{Software and Experimental Setup}
\label{sec:arch}

The Python implementation separates capture, calibration, perception,
preparation, retargeting and export. FastAPI serves offline jobs and the
browser inspection interface; MuJoCo tools reside in a separate image.
Each stage writes a durable episode artifact. The \texttt{rec.npz} hand
extraction, \texttt{cln.npz} prepared motion,
\texttt{object\_models.npz} scene model and \texttt{plan.h5} robot plan
represent different stages and are not interchangeable evidence.
Physical-robot replay is a stub and is not evaluated.

The local apparatus has two Azure Kinect DK cameras on independent USB
controllers. The intended tabletop task is to transfer a blue cube from
a red support to a green target, using a simulated old-generation UR10
with a Robotiq 2F-85 gripper. No physical UR10 is available for this
study. The earlier local cube recordings and their derived artifacts
are not available for verification, so no numerical result from that
campaign is included in this manuscript. In particular, local
capture-synchronisation, stage-yield, object-model, IK, contact-replay
and learned-policy results are \emph{not} claimed. A new, preserved
local cohort would be required to evaluate that end-to-end chain.

\section{External HO-Cap Hand and Object Evaluation}
\label{sec:hocap-results}

The full external experiment covers 28 single-hand HO-Cap pick-and-place
sequences from nine subjects: 24 right-hand and four left-hand, with 29,706
frames from eight calibrated RealSense views. These are not the local Kinect
cube recordings. The source revision and all 114 numerical JSON files are preserved in
the paper's pre-retarget results snapshot with SHA-256 hashes. Capture
and object protocols are documented in the accompanying dated logs.
Both processing queues completed 28/28 scenes. This is a selected task cohort, not an evaluation
on the official HO-Cap held-out split \cite{hocap}.

\subsection{Hand-Joint Accuracy}

Every one of the 623,826 valid reference joint slots has a finite prepared
prediction. Matched-joint-weighted MPJPE is 10.79~mm and PCK@20/50~mm is
92.77/99.85\%. The per-sequence MPJPE range is 9.26--12.85~mm; the
per-sequence P95 range is 17.72--31.90~mm. No pooled P95 is inferred from
these 28 quantiles. The four left-hand scenes have weighted MPJPE
11.60~mm, versus 10.65~mm for 24 right-hand scenes; handedness and subject
are confounded here, so this is descriptive rather than a controlled
difference. The fused-observed stage, before final temporal smoothing,
matches 623,799/623,826 slots (99.996\%) with weighted MPJPE 11.16~mm.
The all-finite smoothed score includes 27 filled slots and should not be
mistaken for direct observation coverage. Table~\ref{tab:hocap-hand-summary}
summarizes the cohort; all 28 episode rows, including P95 and PCK, appear
in Table~\ref{tab:hocap-hand-detail}.

\begin{table}[htbp]
\caption{Eight-view HO-Cap 21-joint accuracy against published MANO fits}
\label{tab:hocap-hand-summary}
\centering\small
\begin{tabular}{lrrrr}
\toprule
Subset and stage & Scenes & Matched / GT & MPJPE (mm) & PCK@20 (\%) \\
\midrule
All, fused observed & 28 & 623,799/623,826 & 11.16 & 91.68 \\
All, smoothed all & 28 & 623,826/623,826 & 10.79 & 92.77 \\
Left, smoothed all & 4 & 86,667/86,667 & 11.60 & 87.56 \\
Right, smoothed all & 24 & 537,159/537,159 & 10.65 & 93.61 \\
\bottomrule
\end{tabular}
\end{table}

\subsection{Named Object Tracks and Outliers}

Each scene contains four annotated objects, giving 112 object--sequence
tracks across 15 object groups. The official HO-Cap object descriptors
\cite{hocap-catalog} are keyed by IDs in Table~\ref{tab:hocap-object-catalog};
for example, \texttt{G18\_2} is a projector remote, not an anonymous
object. All four instances in the \texttt{G18} scenes are the green tea
latte mix, projector remote, right shoe and left shoe. This identity
mapping comes from the authors' \texttt{hocap\_info.yaml}, rather than
appearance-based naming.

The tracker covers 117,916 of 118,824 valid reference object frames
(99.24\%); 97/112 tracks have complete coverage. Across the 112
\emph{track means}, the median and arithmetic mean relative translation
errors are 4.74 and 5.44~mm; the corresponding rotation values are
3.47 and 7.58 degrees. Frame-count-weighted averages of the track means
are 5.60~mm and 7.55 degrees. These are not pooled frame-level P95
values. The median of the 112 track P95s is 10.60~mm translation and
9.25 degrees rotation. First-frame alignment is fitted using the
reference pose, so the numbers measure subsequent relative drift rather
than absolute object-frame accuracy. The supplied instance masks
remove segmentation errors from this test; static intervals remain mixed
with manipulated intervals.

\begin{table}[htbp]
\caption{Selected HO-Cap object failures; mean and P95 are within-track errors}
\label{tab:hocap-object-outliers}
\centering\scriptsize
\begin{tabular}{llrrrrr}
\toprule
Sequence & Object and official descriptor & Cov. (\%) & T mean & T P95 & R mean & R P95 \\
\midrule
S7/162803 & \texttt{G18\_2} projector remote & 98.63 & 12.38 & 20.44 & 112.60 & 175.26 \\
S5/112303 & \texttt{G15\_2} dandelion tea & 100.00 & 17.24 & 21.81 & 109.65 & 128.43 \\
S2/202617 & \texttt{G09\_3} toothpaste box & 100.00 & 5.77 & 13.50 & 74.03 & 132.88 \\
S7/193506 & \texttt{G02\_4} sauce bottle & 99.65 & 2.82 & 10.67 & 56.57 & 169.89 \\
S2/163929 & \texttt{G19\_1} electric screwdriver & 48.80 & 5.75 & 15.41 & 10.55 & 50.95 \\
S2/201942 & \texttt{G06\_1} DVD & 100.00 & 19.39 & 25.24 & 1.90 & 5.63 \\
\bottomrule
\end{tabular}
\end{table}

The worst rotation case is the projector remote in
\texttt{subject\_7/162803}: mean 112.60 degrees and P95 175.26 degrees
despite 98.63\% tracked coverage and 12.38~mm mean translation error.
In another \texttt{G18} scene (\texttt{subject\_1/165502}), the same
object ID has 6.97-degree mean rotation error, and the other three
objects in the failing scene have 1.10--1.34-degree means. Thus the
anomaly is sequence/object-specific rather than a global frame rotation.
A visual review showed plausible late object placement under an alternate
orientation, but did not establish a unique root cause; the outlier remains
in all aggregates. The dandelion tea (109.65 degrees) and toothpaste box
(74.03 degrees in one scene, versus 13.82 and 6.89 degrees in two others)
show that this is not a single isolated failure. A low translation error
also does not certify orientation, as the sauce bottle demonstrates.
Conversely, the DVD has the highest mean translation error, 19.39~mm,
with only 1.90-degree mean rotation error. The electric screwdriver is a
coverage failure at 48.80\%, so its conditional error on visible frames
must not stand in for whole-sequence accuracy. Tables
\ref{tab:hocap-object-detail} and \ref{tab:hocap-object-catalog} retain
all 112 object scores and official identities for audit.

\subsection{Historical Two-Sequence Camera-Pair Pilot}

Before the full eight-view run, two subject-1 sequences were tested for
120 frames each with a pair selected after screening all 28 pairs on those
same frames. Table~\ref{tab:hocap-hand} is retained as a historical
\emph{in-sample} development comparison; its original local artifacts
were lost. The full-sequence eight-view results above supersede it as the
reproducible hand/object reference cohort, but do not turn the selected
pair into a held-out baseline.

\begin{table}[htbp]
\caption{Historical HO-Cap Hand-Joint Pilot: Selected Pair Versus Eight Views}
\label{tab:hocap-hand}
\centering\small
\begin{tabular}{llrrrr}
\toprule
Sequence & Views & Coverage & MPJPE (mm) & P95 (mm) & PCK@20 (\%) \\
\midrule
\texttt{170650} & selected pair & 2,520/2,520 & 11.33 & 22.71 & 87.70 \\
\texttt{170650} & all eight & 2,520/2,520 & 11.34 & 22.17 & 87.50 \\
\texttt{170105} & selected pair & 2,520/2,520 & 11.67 & 22.80 & 89.80 \\
\texttt{170105} & all eight & 2,520/2,520 & 10.50 & 20.15 & 94.88 \\
\bottomrule
\end{tabular}
\end{table}

The selected pair also failed one colour--depth consistency check. No
HO-Cap score establishes the accuracy of the local Kinect calibration or
physical success of the proposed robot motion.


\section{Downstream Evidence Boundary}
\label{sec:downstream-boundary}

The archived HO-Cap experiment evaluates hand reconstruction and
first-frame-aligned object tracking on a separate RealSense capture
system. It does not establish the accuracy of the local Kinect rig or
produce a validated robot demonstration. Retarget plans, simulator
contacts, task success and policy learning require their own retained
inputs, protocols and numerical artifacts before they can be reported
as results. This chapter therefore ends at the external perception
evaluation rather than extrapolating an end-to-end success rate.
